\chapter{Бін \texorpdfstring{$\neq$}{≠} плазма, але корисний контроль}\label{ch:g03}

Це розділ про найповчальніше спростування проєкту. Ми висунули красиву гіпотезу про
спільну математику надпровідників і плазми, перевірили її --- і вона впала. Але те, що
лишилося після падіння, дало точний числовий результат, якого в літературі, наскільки
нам відомо, немає: середнє двох розв'язків нелінійної задачі не є розв'язком, і метрика
$L^2$ цього не бачить. Ланцюжок: R2~→ D8~→ E4/E5~→ C7 (рис.~\ref{fig:g00-chains}).

\section{Гіпотеза, яка впала: R2}\label{sec:g03-r2}

Пропозиція (2026-09-22, \texttt{Our try/02-field-map/D-superconductivity.md}) звучала так:
критичний стан Біна в надпровіднику II роду і задача Ґреда--Шафранова з вільною межею
мають спільну структуру «вільна межа~/ варіаційна нерівність», тож той самий чисельний
апарат (Deep Ritz, PINN, варіаційні розв'язувачі) застосовний до обох.

\begin{outofcorpus}[критичний стан Біна (плита, стале $J_c$)]
У моделі Біна густина струму в надпровіднику або дорівнює критичній, $|J|=J_c$, або
нульова. Для плити в паралельному полі $H_a$ профіль поля всередині
\[
H(x)=\max\bigl(H_a-J_c x,\;0\bigr),
\]
фронт проникнення $x_p=H_a/J_c$. Задача --- варіаційна нерівність з \emph{опуклою}
множиною обмежень $K=\{\,|{\pa H}/{\pa x}|\le J_c\,\}$; вільна межа --- це межа
\emph{активної множини}, де обмеження насичене.
\end{outofcorpus}

Перевірка літератури дала три факти, які вирішують справу \cite{RegEst,RegDec}:
\begin{enumerate}
\item \textbf{Опуклість.} Для сталого $J_c$ Прігожин \cite{Prigozhin1996} (Теорема~2)
  переписує задачу Біна як $\dd h/\dd t+\pa I_K(h)\ni f$, де $\pa I_K$ --- субдиференціал
  індикатора опуклої множини, максимально монотонний оператор. Звідси існування й
  \emph{єдиність}. Задача плазми --- напівлінійна нелінійна задача на власні значення з
  \emph{неопуклим} функціоналом; розв'язок у загальному випадку неєдиний \cite{Schaeffer1977,Bartolucci2021}.
\item \textbf{Перешкода проти лінії рівня.} У Біна вільна межа --- межа активної множини
  обмеження-нерівності з множником Лагранжа $\rho\ge0$. У плазми вона --- лінія рівня
  розв'язку $\pa\{v>0\}$: жодної перешкоди, жодного множника при нерівності.
\item \textbf{Нульове перехресне цитування за 50 років:} дві зрілі спільноти одна одну не
  цитують.
\end{enumerate}

\begin{established}{E24 (прочитано)}
Критичний стан Біна опуклий і однозначно розв'язний; задача плазми неопукла й
багатозначна (Prigozhin 1996, Тм.~2; Bartolucci та ін. 2021; Schaeffer 1977) \cite{RegEst}.
\end{established}
\begin{refuted}{R2}
«Критичний стан Біна і вільна межа плазми --- одна математика». Спростовано: опукла
проти неопуклої; активна множина проти лінії рівня; нульове перехресне цитування
\cite{RegEst}. Важливе застереження з огляду: PINN застосовні до обох задач --- але вони
застосовні майже до будь-якого PDE, і спільна застосовність універсального методу не є
доказом спільної структури.
\end{refuted}

\paragraph{Що врятували: D8.} Питання «чи це одна задача» замінили іншим: \emph{чи
відрізняється режим відмови методу між опуклою задачею з вільною межею і неопуклою?}
Бін стає \textbf{опуклим контролем}, проти якого видно наслідки неопуклості. Рішення D8
закрило надпровідність як окрему тему (2026-09-22) і злило її з глибоким розбором D2
\cite{RegDec}. Реєстр гіпотез прямо називає R2 одним із двох спростувань, які дали
\emph{кращий} результат, ніж містили.

\section{Навіщо неєдиність: аргумент проти \texorpdfstring{$L^2$}{L2}}\label{sec:g03-why}

\begin{established}{E21 (прочитано)}
Розв'язок рівняння Ґреда--Шафранова може бути неєдиним на реальній геометрії: Ham і
Farrell 2024 \cite{HamFarrell2024}; Pentland та ін. 2025 \cite{Pentland2025} (MAST-U,
дефляційна континуація) \cite{RegEst}.
\end{established}
\noindent Робота по MAST-U не містить машинного навчання взагалі, але містить потрібний
нам факт: коли підгонку обмежують лише магнітні виміри, дві дуже різні форми профілю
можуть давати однаковий потік на осі. Тобто відображення «діагностики~→ $\psi$», яке
вчить кожна ML-модель реконструкції, може бути не функцією, а \emph{багатозначним
відношенням}.

Регресор, що мінімізує $\mathbb{E}\|f(x)-y\|^2$, збігається до умовного середнього
$\mathbb{E}[y\,|\,x]$. Якщо за тих самих $x$ існують дві гілки $y_1,y_2$, модель
повертає щось на кшталт $\tfrac12(y_1+y_2)$. Для лінійного рівняння це теж розв'язок.
Для нелінійного --- ні, і ось чому.

\begin{derivation}[нев'язка середнього двох гілок]
Нехай $u_1,u_2$ розв'язують $u''+\lambda\ee^{u}=0$. Для
$\bar u=\tfrac12(u_1+u_2)$ маємо $\bar u''=-\tfrac{\lambda}{2}(\ee^{u_1}+\ee^{u_2})$, тож
\begin{equation}\label{eq:g03-jensen}
r[\bar u]\equiv\bar u''+\lambda\ee^{\bar u}
=\lambda\Bigl(\ee^{\bar u}-\tfrac12\bigl(\ee^{u_1}+\ee^{u_2}\bigr)\Bigr)\le0 .
\end{equation}
Нерівність --- це нерівність Єнсена для опуклої $\ee^{u}$; рівність лише при $u_1=u_2$.
Отже, середнє порушує рівняння \emph{всюди}, де гілки різні, і тим сильніше, чим далі
вони одна від одної. А в $L^2$ середнє --- найкраща можлива відповідь \emph{за
означенням}: жодне збільшення ваги $L^2$-члена цього не виявить. Для рівняння
Ґреда--Шафранова $\GSop\psi=-\mu_0R^2p'(\psi)-FF'(\psi)$ механізм той самий: права
частина нелінійна за $\psi$ (\cite{Manual}, розд.~3).
\end{derivation}

\section{Сурогат механізму: Братý проти Біна}\label{sec:g03-bratu}

Для точного експерименту потрібна неопукла задача, чиї гілки відомі в замкненій формі.
Такою є одновимірна задача Братý (Гельфанда)
\begin{equation}\label{eq:g03-bratu}
-u''=\lambda\ee^{u},\qquad u(0)=u(1)=0 .
\end{equation}
Bartolucci та ін. \cite{Bartolucci2021} відносять задачу плазми з вільною межею саме до
сімейства Гельфанда~/ середньопольових вихорів~/ напівлінійних біфуркацій, тож це
правильний родич, хоча й не сама GS.

\begin{outofcorpus}[гілки Братý і точка складки]
Розв'язок \eqref{eq:g03-bratu} має вигляд
$u(x)=-2\ln\bigl[\cosh(s\theta/2)/\cosh(\theta/4)\bigr]$, $s=x-\tfrac12$, де $\theta$
задовольняє $\theta=\sqrt{2\lambda}\cosh(\theta/4)$. Функція $\theta/\cosh(\theta/4)$ спершу
зростає, потім спадає, тож при $\lambda<\lambda^{*}$ рівняння має рівно два корені (дві
гілки), при $\lambda=\lambda^{*}$ --- один, далі жодного. Це \emph{складка} (fold). У точці
складки похідна $\theta/\cosh(\theta/4)$ нульова: $\cosh(\theta_c/4)=(\theta_c/4)\sinh(\theta_c/4)$,
$\lambda^{*}=\theta_c^2/\bigl(2\cosh^2(\theta_c/4)\bigr)$. Енергія
$J[u]=\int_0^1\bigl(\tfrac12u'^2-\lambda\ee^{u}\bigr)\dd x$ знаконевизначена й необмежена
знизу (візьміть $u=t\varphi$, $t\to\infty$); нижня гілка --- її локальний мінімум,
верхня --- сідлова точка («гірський перевал»).
\end{outofcorpus}

Контроль --- плита Біна зі сталим $J_c$: один розв'язок на кожне $H_a$, тож «умовне
середнє по множині розв'язків» просто збігається з розв'язком.

\subsection{Код: точні гілки}

\code{ourtry/03-deep-dives/D2-gs-nonuniqueness/branch_mean.py}{61}{83}{(корені $\theta$, гілки, точка складки)}

\begin{itemize}
\item \emph{Рядки 61--70, \texttt{theta\_roots}:} шукаємо обидва корені
  $g(\theta)=\theta-\sqrt{2\lambda}\cosh(\theta/4)$. Замість «вгадування» дужок --- густа
  сітка з 200\,000 точок на $[0,60]$, зміни знака $g$ і уточнення кожного кореня методом
  Брента (\texttt{brentq}) до $10^{-14}$. При $\lambda>\lambda^{*}$ змін знака немає ---
  список порожній.
\item \emph{Рядки 73--75, \texttt{bratu\_u}:} точна формула гілки для даного $\theta$.
\item \emph{Рядки 78--83, \texttt{lam\_star}:} точка складки з умови
  $\cosh(\theta/4)=(\theta/4)\sinh(\theta/4)$ --- знову \texttt{brentq}, потім
  $\lambda^{*}=\theta_c^2/(2\cosh^2(\theta_c/4))$. Це і є валідація реалізації (E5).
\end{itemize}

\subsection{Код: нев'язки}

\code{ourtry/03-deep-dives/D2-gs-nonuniqueness/branch_mean.py}{86}{104}{(профіль Біна, друга похідна, нев'язка Братý)}

\begin{itemize}
\item \emph{Рядки 87--89, \texttt{bean\_H}:} точний профіль Біна $\max(H_a-J_cx,0)$.
\item \emph{Рядки 93--97, \texttt{d2}:} друга похідна четвертого порядку точності
  (п'ятиточковий шаблон) на сітці з $N=20\,001$ вузла; крайні два вузли --- \texttt{NaN}.
\item \emph{Рядки 100--104, \texttt{rel\_residual\_bratu}:} відносна нев'язка
  $\|u''+\lambda\ee^{u}\|/\|\lambda\ee^{u}\|$ --- у частках норми джерела, з відступом
  200 вузлів від країв. Значення 1 означає, що нев'язка такого самого розміру, як сам
  член рівняння.
\end{itemize}

\code{ourtry/03-deep-dives/D2-gs-nonuniqueness/branch_mean.py}{107}{117}{(нев'язка Біна)}

\noindent \emph{Рядки 107--117:} для Біна перевіряємо $\dd H/\dd x+J_c=0$ лише в
проникненій області ($H>10^{-6}$), відступивши по 5 вузлів від країв і від злому на фронті;
сама вільна межа виключена. Нормування інше, ніж у Братý (на $J_c\sqrt n$), тож числа двох
задач порівнюємо за порядком, а не буквально.

\subsection{Код: експеримент}

\code{ourtry/03-deep-dives/D2-gs-nonuniqueness/branch_mean.py}{121}{145}{(неопуклий тест)}

\begin{itemize}
\item \emph{Рядки 122--124:} друкуємо $\lambda^{*}$ поруч із літературним 3,513830719.
\item \emph{Рядки 130--139:} для $\lambda\in\{0{,}5;\,1;\,2;\,3;\,3{,}4;\,3{,}5\}$ беремо обидві
  гілки, їхнє середнє \texttt{u\_mean} (те, що поверне $L^2$-регресор за рівноймовірних
  гілок) і три нев'язки.
\item \emph{Рядки 140--145:} \texttt{L2 sep} --- середньоквадратична відстань між гілками;
  усе пишеться в \texttt{rows}.
\end{itemize}

\code{ourtry/03-deep-dives/D2-gs-nonuniqueness/branch_mean.py}{147}{171}{(опуклий контроль і підсумок)}

\begin{itemize}
\item \emph{Рядки 149--158:} для п'яти значень $H_a$ «середнє по множині розв'язків» ---
  копія єдиного розв'язку (рядок 152).
\item \emph{Рядки 164--171:} головне порівняння: середня нев'язка справжньої гілки,
  середня нев'язка $L^2$-середнього (усереднено по шести $\lambda$) і те саме для Біна.
\end{itemize}

\section{Що вийшло}\label{sec:g03-results}

Прогін 2026-09-22 (\texttt{results.json}); функції \texttt{lam\_star}, \texttt{theta\_roots}
і нев'язки при $\lambda=0{,}5$ і $3{,}5$ ми 2026-09-29 перерахували імпортом модуля (без
запуску \texttt{main}, щоб не перезаписати файл) --- збіг до останньої цифри.

\begin{center}\small
\begin{tabular}{@{}ccccccc@{}}
\toprule
$\lambda$ & $\theta_{\mathrm{lo}}$ & $\theta_{\mathrm{hi}}$ & нев'язка (нижня) & нев'язка (верхня)
  & нев'язка середнього & відстань гілок \\
\midrule
0,5 & 1,0336 & 13,0382 & $3{,}40\cdot10^{-7}$ & $1{,}10\cdot10^{-8}$ & 4,2609 & 3,3624 \\
1,0 & 1,5172 & 10,9387 & $1{,}64\cdot10^{-7}$ & $1{,}26\cdot10^{-8}$ & 2,1018 & 2,6545 \\
2,0 & 2,3576 & 8,5072 & $7{,}22\cdot10^{-8}$ & $1{,}41\cdot10^{-8}$ & 0,7249 & 1,7514 \\
3,0 & 3,3735 & 6,5766 & $3{,}95\cdot10^{-8}$ & $1{,}63\cdot10^{-8}$ & 0,1745 & 0,9207 \\
3,4 & 4,1084 & 5,5623 & $2{,}94\cdot10^{-8}$ & $1{,}88\cdot10^{-8}$ & 0,0347 & 0,4193 \\
3,5 & 4,5519 & 5,0543 & $2{,}52\cdot10^{-8}$ & $2{,}37\cdot10^{-8}$ & 0,0041 & 0,1450 \\
\bottomrule
\end{tabular}
\end{center}

\begin{established}{E4}
$L^2$-середнє двох гілок не є розв'язком, і $L^2$ цього не бачить: відносна нев'язка
середнього \textbf{1,22} (середнє по шести $\lambda$) проти \textbf{$6{,}4\cdot10^{-8}$} у
справжньої гілки, тобто в $1{,}9\cdot10^{7}$ разів гірше; опуклий контроль Біна ---
\textbf{$2{,}5\cdot10^{-13}$} \cite{RegEst}.
\end{established}
\begin{established}{E5}
Реалізація валідна: точка складки Братý $\lambda^{*}=3{,}513830719$ збігається з
літературним значенням до останньої цифри ($\theta_c=4{,}7987$) \cite{RegEst}.
\end{established}

\begin{figure}[ht]
\centering
\includegraphics[width=\textwidth]{g03_branch_mean.jpg}
\caption{(a)~Дві точні гілки Братý при $\lambda=1$ і їхнє $L^2$-середнє; (b)~нев'язка
$u''+\lambda\ee^{u}$: у гілок нуль, у середнього --- від'ємна всюди, як
у~\eqref{eq:g03-jensen}; (c)~відносна нев'язка середнього і гілки проти $\lambda$, опуклий
контроль Біна і відстань між гілками (права вісь). Власний розрахунок, код:
\texttt{Our try/03-deep-dives/D2-gs-nonuniqueness/\{branch\_mean,plots\}.py}.}
\label{fig:g03-branch}
\end{figure}

Як це читати. Нев'язка 1,22 означає, що «розв'язок» порушує рівняння на величину порядку
самого джерела; при $\lambda=0{,}5$ --- учетверо більше. Структура залежності правильна:
при $\lambda\to\lambda^{*}$ гілки зливаються (відстань $3{,}36\to0{,}145$), і нев'язка
середнього падає до 0,0041, як і вимагає~\eqref{eq:g03-jensen}. Інформацію знищують не дані
і не потужність мережі: кожна окрема гілка має нев'язку на рівні числового шуму. Її знищує
\emph{цільова функція}. Правильні інструменти для багатозначних цілей відомі --- суміші
розподілів (MDN), умовні генеративні моделі, оборотні мережі.

\paragraph{Межі результату --- читати обов'язково.}
\begin{itemize}
\item Це \emph{не} Ґред--Шафранов. Перенесення висновку на GS спирається на цитати
  \cite{HamFarrell2024,Pentland2025,Bartolucci2021}, а не на цей розрахунок. Названий
  наступний крок --- дефляційна континуація поверх \texttt{FreeGSNKE} на геометрії MAST-U
  за рецептом Pentland та ін. --- не зроблено.
\item Контроль Біна в коді --- тавтологія: множина розв'язків одноелементна \emph{за
  побудовою} (рядок 152), тож рівність нев'язок розв'язку й «середнього» гарантована, а
  $2{,}5\cdot10^{-13}$ --- лише точність чисельної похідної кусково-лінійної функції. Уся
  змістовна частина контролю --- теорема Прігожина (E24); код її ілюструє, а не перевіряє.
\item Модель Біна--Кіма ($J_c$ залежить від поля) квазіваріаційна; для неї жодного
  аналогічного виміру немає.
\end{itemize}

\begin{itemcard}{Q3}{відкрите питання}
Чи неєдиність розв'язку GS (Ham--Farrell 2024, Pentland 2025) зустрічається в робочих
режимах, чи лише в екзотичних? Від цього залежить, чи розбір D2 --- реальна задача, чи
математичний курйоз. Адресат --- фізики плазми; статус --- відкрите \cite{RegQ}.
\end{itemcard}

\section{Наступний крок: C7 і Deep Ritz}\label{sec:g03-c7}

\begin{conjecture}{C7}
Deep Ritz на задачі плазми буде seed-нестабільним (обиратиме різні гілки), а на Біні ---
стабільним. \emph{Спростує:} обидва стабільні --- неопуклість не проявляється на доступних
масштабах. \emph{Вартість:} $\approx$ день, потрібні обидва розв'язувачі \cite{RegConj}.
\end{conjecture}

Метод Deep Ritz \cite{EYu2018} подає розв'язок нейромережею $u_\theta$ і мінімізує
енергію варіаційної задачі, оцінену Монте-Карло, плюс штраф за крайову умову. Приклад у
\texttt{Tokamak-GS-solver} --- загальний (десятивимірне рівняння Лапласа з точним
розв'язком $x_1x_2+x_3x_4+\ldots$), не GS і не Братý, але вся механіка на місці.

\code{gs/src/examples/deep_ritznet.py}{238}{259}{(цикл навчання Deep Ritz)}

\begin{itemize}
\item \emph{Рядки 241--243:} точки колокації всередині кулі (\texttt{data1}, з
  \texttt{requires\_grad}, бо знадобиться $\grad u$) і на межі (\texttt{data2}).
\item \emph{Рядки 246--250:} прямий прохід і $\grad u_\theta$ через
  \texttt{torch.autograd.grad} з \texttt{create\_graph=True} --- похідна має бути
  диференційовною за вагами.
\item \emph{Рядки 252--253:} енергія Рітца $\mathbb{E}\bigl[\tfrac12|\grad u|^2-fu\bigr]$ для
  $-\Delta u=f$ (тут $f=0$, функція \texttt{ffun}).
\item \emph{Рядки 256--259:} штраф $\beta\cdot$\,площа$\cdot\mathbb{E}(u-g)^2$ на межі
  ($\beta=500$) і повна втрата.
\end{itemize}
Зерно задано в рядку 327 (\texttt{torch.manual\_seed(21)}); саме його треба змінювати, щоб
міряти «дисперсію гілки по зернах».

\paragraph{Зауваження путівника до постановки C7.} Для Братý заміна \texttt{fTerm*output1}
на $\lambda\ee^{u}$ дає енергію $J[u]$, яка необмежена знизу (див. рамку вище). Чиста
мінімізація може знайти лише \emph{нижню} гілку (локальний мінімум) або «втекти» на
нескінченність; верхня гілка --- сідло, і градієнтний спуск її не обирає. Тому на Братý
C7 швидше проявиться як «збіжність до нижньої гілки або розбіжність залежно від зерна»,
ніж як вибір між двома гілками. Для справжньої задачі плазми Bartolucci та ін. мінімізують
на неопуклому многовиді зі сталою повного струму, і картина може бути іншою. Це
міркування, а не вимір; реєстр його не містить.

\begin{practicum}[C7: seed-стабільність Deep Ritz]
\begin{enumerate}
\item \textbf{Відтворення D2} (безпечно, у копії):
  скопіюйте \texttt{branch\_mean.py} і \texttt{plots.py} у власну теку й запустіть
  інтерпретатором стартового набору (розд.~\ref{sec:g00-run}); \texttt{results.json} має
  збігтися з наведеною таблицею.
\item \textbf{Контроль.} Одновимірна плита Біна з регуляризацією степеневим законом
  $E\propto(J/J_c)^n$ (опукла енергія). Deep Ritz на 10 зернах; міряйте положення фронту
  $x_p$ проти точного $H_a/J_c$. Очікування за E24: розкид мізерний.
\item \textbf{Тест.} Братý при $\lambda=1$: Deep Ritz на енергії $J[u]$ з 10 зернами. Для
  кожного зерна класифікуйте результат: нижня гілка ($u(\tfrac12)\approx0{,}141$),
  верхня ($u(\tfrac12)\approx4{,}091$; обидва значення --- з \texttt{bratu\_u}), розбіжність. Потім спробуйте варіант, що може знайти
  верхню гілку (наприклад, PINN на сильній нев'язці, де обидві гілки --- нулі втрати).
\item \textbf{Звіт у стилі реєстру:} що виміряно, чи спростовано C7 і в який бік; якщо
  результат --- «на Братý Deep Ritz бачить лише нижню гілку», це теж відповідь, і її треба
  записати, а не відкинути.
\end{enumerate}
\end{practicum}
